%! TEX program = xelatex
%! TEX root = ../main.tex
%! TEX encoding = utf-8

%%%%%%%%%%%%%%%%%%%%%%%%%%%%%%%%%%%%%%%%%%%%%%%%%%%%%%%%%%%%%%%%%%%%%%
%
%  哈尔滨工程大学学位论文 XeLaTeX 模版 —— 正文文件 chap02.tex
%
%  版本：1.0.0
%  最后更新：
%  修改者：Leo LiWenhui lwh@hrbeu.edu.cn
%  修订者：
%  编译环境1：Ubuntu 12.04 + TeXLive 2013/2014
%  编译环境2：Windows 7/8  + TeXLive 2013/2014
%
%%%%%%%%%%%%%%%%%%%%%%%%%%%%%%%%%%%%%%%%%%%%%%%%%%%%%%%%%%%%%%%%%%%%%

\chapter{浅海水声信道}[Shallow-Water Acoustic Channel]
\label{chap02}

\section{浅海水声传播特点}[Propagation Characteristics]

浅海环境通常指水深远小于传播距离且上下边界对传播过程具有显著约束作用的海洋区域。在该环境中，声波传播受到海面起伏、海底反射、温盐深结构以及平台运动等因素共同影响，传播路径通常由直达径、海面反射径、海底反射径及多次反射径组成，接收端易形成明显的时延扩展和相位起伏\cite{hovemRayTraceModeling2013,nosalMeasuredModeledHighfrequency2010}。

对声学释放器而言，浅海信道的突出特征是多径密集和环境时变。若符号持续时间短于主要多径时延扩展，接收端将出现严重码间干扰；若多个反射径与直达径能量接近，则首达径检测和同步捕获都会变得困难。与此同时，海浪、船体漂移和目标缓慢运动还会引入不可忽略的多普勒频偏，进而影响跳频同步与频点判决。

\section{噪声与干扰特性}[Noise and Interference]

浅海背景噪声一般由风浪噪声、航运噪声、生物噪声及设备自噪声叠加而成，其功率谱密度随频率、海况和作业海域变化而变化。声学释放器工作频段多处于中低频至中频范围，既要避免低频强噪声区，又要兼顾换能器效率和传播损耗。除背景噪声外，邻频设备辐射、主动声纳脉冲和局部机电系统谐波还可能表现为持续窄带干扰或突发干扰，这类干扰对单载频通信体制影响尤为明显\cite{freitagAnalysisChannelEffects2001,SuMengTiaoPinKuoPinTongXinXiTongDeFangZhenYanJiu2024}。

FH-MFSK 通过在时频域分散能量，可显著降低部分频带干扰造成的链路失效风险。当局部频带受到压制时，仅部分跳频单元受损，因此系统仍有机会通过多跳综合判决恢复正确指令。这也是选择 FH-MFSK 作为声学释放器通信体制的重要原因之一。

\section{信道时变性及多普勒效应}[Time Variability and Doppler]

若发射端与接收端存在相对径向速度 $v$，中心频率为 $f_c$，介质平均声速为 $c$，则一阶近似下的多普勒频移可写为

\begin{equation}
  \Delta f \approx \frac{v}{c} f_c.
  \label{eq:doppler}
\end{equation}

虽然平台速度通常远小于声速，但当载频较高、频点间隔较小时，多普勒频移依然可能导致接收能量偏离理想滤波器带宽中心，进而降低非相干检测输出。若存在海面动态反射与多径差异运动，还会形成不同路径的频移不一致现象，使同步与判决问题进一步复杂化\cite{TongFengYiZhongQianHaiXinDaoTiaoPinShuiShengDiaoZhiJieDiaoXiTong2007}。

因此，在工程设计中应综合考虑以下因素：其一，频点间隔需大于最大频偏与接收滤波器主瓣宽度之和；其二，单跳持续时间不宜过长，以降低时变信道对同一符号能量积累的破坏；其三，帧结构中需要设置可用于粗同步和频偏估计的前导段。

\section{浅海信道对 FH-MFSK 参数设计的约束}[Design Constraints]

设系统总可用带宽为 $B$，每跳可选频点数为 $M$，单个子频点有效占用带宽为 $B_s$，保护带宽为 $B_g$，则近似有

\begin{equation}
  B \geq M(B_s + B_g).
  \label{eq:bandwidth_constraint}
\end{equation}

符号持续时间 $T_s$ 与频率分辨率互相制约。若 $T_s$ 过短，则接收积分增益不足；若 $T_s$ 过长，则更易受到多径扩展和时变衰落影响。结合浅海信道特点，FH-MFSK 参数设计通常需在以下目标之间取得平衡：

\begin{enumerate}
  \item 足够的频点间隔，以保证频率正交性和抗频偏能力；
  \item 适中的跳速，以兼顾抗干扰能力与同步复杂度；
  \item 有限的帧长，以缩短一次完整指令传输时间并降低误码积累；
  \item 合理的保护间隔或相关门限，以提高多径环境下的判决稳定性。
\end{enumerate}

\section{面向测距的传播时延建模}[Propagation Delay Model for Ranging]

测距算法的基础是传播时延建模。若忽略复杂折射效应并采用等效平均声速 $c_0$，目标斜距 $R$ 与传播时延 $\tau$ 满足

\begin{equation}
  R = c_0 \tau.
  \label{eq:simple_range}
\end{equation}

然而在真实浅海环境中，声速会随温度、盐度和压力变化，声线也可能发生弯折，导致式\eqref{eq:simple_range}仅适合作为一阶近似。若测距精度要求提高，则需进一步引入声速剖面修正项、应答机固定延时项和首达径估计误差项\cite{sunRelationshipPropagationTime2019,bergerStratificationEffectCompensation2008}。

\section*{本章小结}[Summary]

本章分析了浅海水声信道的多径、噪声、时变和多普勒等关键特征，并从带宽、跳速、频点间隔和时延建模等角度讨论了其对 FH-MFSK 声学释放器通信与测距设计的影响，为后续算法研究奠定了信道基础。

（1）使用 xelatex

\begin{lstlisting}
    $ xelatex spine
    $ xelatex a3cover
\end{lstlisting}

（2）或者使用 latexmk

\begin{lstlisting}
    $ latexmk spine
    $ latexmk a3cover
\end{lstlisting}

在 \texttt{Windows} 系统下也可以直接使用 \texttt{build\_all.bat} 批处理生成 \texttt{A3} 封面。

\section{字库安装}[Font Install]

可以通过模板选项选择~\texttt{fontset}~使用Windows系统字体或者Adobe字体。

Windows系统下本模板默认是使用Windows字体字库，可以不用设置字体选项，如果使用Adobe字体字库，需要设置~\texttt{fontset=adobe}~。
在使用此模板撰写论文前，应该确保相应的字库已经安装，并且最好是包含宋体、黑体、楷体和仿宋的完整套装。

在 Windows 操作系统下，只要把字库文件复制到 Windows 的 Fonts 文件夹下即可，
而对于 Linux 系统，可通过右键点击字库文件然后选择【安装字库】菜单选项进行安装。
Linux 对于系统新安装的字库，需要使用命令~\texttt{sudo fc-cache -fsv}~刷新缓存后才可以使用。

\section*{本章小结}[Brief Summary]
\LaTeX{}~工作环境安装与配置简介。
